\documentclass[12pt,letterpaper]{article}

\usepackage[spanish]{babel}
\usepackage[utf8]{inputenc}
\usepackage[T1]{fontenc}
\usepackage{graphicx}
\usepackage{fancyhdr}
\usepackage{geometry}
\usepackage{amsmath, amssymb}

\geometry{margin=2.5cm}
\setlength{\headheight}{52pt}


\pagestyle{fancy}
\fancyhf{}

\fancyhead[L]{
    \includegraphics[height=1.5cm]{ur_image.png}
}

\fancyhead[R]{
\begin{tabular}{r}
\textbf{Monitoria Álgebra Lineal} \\
Gabriel Fernando Salcedo Zamora \\
\end{tabular}
}

\renewcommand{\headrulewidth}{0.4pt}

\begin{document}

\section*{Taller: Espacios Fundamentales y Rango de una Matriz}

\begin{enumerate}

    %----------------------------------------------------
    % Espacio fila, espacio columna, espacio nulo
    %----------------------------------------------------

    \item Para cada una de las siguientes matrices, encuentre una base para el \textbf{espacio fila}, el \textbf{espacio columna} y el \textbf{espacio nulo}. Verifique además el Teorema Rango-Nulidad.

    \begin{enumerate}
        \item $A = \begin{pmatrix} 1 & 2 & 3 \\ 4 & 5 & 6 \\ 7 & 8 & 9 \end{pmatrix}$

        \item $B = \begin{pmatrix} 1 & 0 & -2 & 1 \\ 0 & 1 & 3 & -2 \\ 2 & -1 & -7 & 4 \end{pmatrix}$

        \item $C = \begin{pmatrix} 2 & -1 & 3 \\ 4 & -2 & 6 \end{pmatrix}$
    \end{enumerate}

    %----------------------------------------------------
    % Rango de una matriz
    %----------------------------------------------------

    \item Determine el rango de cada una de las siguientes matrices. Justifique su respuesta usando la forma escalonada reducida.

    \begin{enumerate}
        \item $\begin{pmatrix} 1 & 3 & -2 \\ 2 & -1 & 4 \\ -1 & 4 & -6 \end{pmatrix}$

        \item $\begin{pmatrix} 0 & 1 & 2 & 3 \\ 1 & 0 & -1 & 2 \\ 2 & 1 & 0 & 7 \\ 3 & 2 & 1 & 12 \end{pmatrix}$

        \item $\begin{pmatrix} a & 0 & 0 \\ 0 & b & 0 \\ 0 & 0 & c \end{pmatrix}$, donde $a, b, c \in \mathbb{R}$. Analice todos los casos posibles.
    \end{enumerate}

    %----------------------------------------------------
    % Sistemas lineales y espacio nulo
    %----------------------------------------------------

    \item Considere el sistema de ecuaciones $A\mathbf{x} = \mathbf{0}$, donde
    \[
    A = \begin{pmatrix} 1 & -1 & 2 & 1 \\ 2 & -2 & 4 & 2 \\ 3 & 1 & -1 & 0 \\ 1 & 3 & -5 & -1 \end{pmatrix}.
    \]
    \begin{enumerate}
        \item Encuentre la forma escalonada reducida de $A$.
        \item Identifique las variables libres y las variables básicas (pivote).
        \item Encuentre una base para el espacio nulo de $A$ y determine su dimensión (nulidad).
        \item Calcule el rango de $A$ y verifique que $\text{rango}(A) + \text{nulidad}(A) = n$, donde $n$ es el número de columnas.
    \end{enumerate}

    %----------------------------------------------------
    % Consistencia y espacio columna
    %----------------------------------------------------
\newpage
    \item Sea $A = \begin{pmatrix} 1 & 2 & -1 \\ 3 & 1 & 2 \\ -1 & 3 & -4 \end{pmatrix}$.
    Para cada uno de los siguientes vectores $\mathbf{b}$, determine si el sistema $A\mathbf{x} = \mathbf{b}$ es consistente, verificando si $\mathbf{b}$ pertenece al espacio columna de $A$.

    \begin{enumerate}
        \item $\mathbf{b} = (1,\, 2,\, 3)^T$
        \item $\mathbf{b} = (5,\, 8,\, 2)^T$
        \item $\mathbf{b} = (0,\, 0,\, 0)^T$
    \end{enumerate}

    %----------------------------------------------------
    % Espacio nulo izquierdo
    %----------------------------------------------------

    \item Para la matriz
    \[
    M = \begin{pmatrix} 1 & 2 & 3 \\ 0 & 1 & 4 \\ 0 & 0 & 0 \end{pmatrix},
    \]
    \begin{enumerate}
        \item Encuentre una base para cada uno de los cuatro espacios fundamentales de $M$: espacio fila, espacio columna, espacio nulo y espacio nulo izquierdo (espacio nulo de $M^T$).
        \item Organice los resultados en una tabla que muestre el espacio, su base y su dimensión.
        \item Verifique que las dimensiones satisfacen el Teorema Rango-Nulidad tanto para $M$ como para $M^T$.
    \end{enumerate}

\end{enumerate}

%\section*{Solución}
\newpage

\section*{Solución}

\begin{enumerate}

%-------------------------------------------------
% Solución 1: Espacios fundamentales
%-------------------------------------------------

\item Espacios fundamentales y verificación del Teorema Rango-Nulidad.

\begin{enumerate}

    \item Para $A = \begin{pmatrix} 1 & 2 & 3 \\ 4 & 5 & 6 \\ 7 & 8 & 9 \end{pmatrix}$, aplicamos eliminación de Gauss-Jordan:
    \[
    A \xrightarrow{F_2-4F_1,\; F_3-7F_1}
    \begin{pmatrix} 1 & 2 & 3 \\ 0 & -3 & -6 \\ 0 & -6 & -12 \end{pmatrix}
    \xrightarrow{F_3-2F_2}
    \begin{pmatrix} 1 & 2 & 3 \\ 0 & -3 & -6 \\ 0 & 0 & 0 \end{pmatrix}
    \xrightarrow{-\frac{1}{3}F_2}
    \begin{pmatrix} 1 & 0 & -1 \\ 0 & 1 & 2 \\ 0 & 0 & 0 \end{pmatrix}
    \]

    \textbf{Espacio fila:} Las filas no nulas de la RREF forman una base:
    \[
    \mathcal{B}_{\text{fila}} = \{(1,0,-1),\; (0,1,2)\}
    \]

    \textbf{Espacio columna:} Las columnas de $A$ correspondientes a las columnas pivote (columnas 1 y 2) de la RREF:
    \[
    \mathcal{B}_{\text{col}} = \left\{\begin{pmatrix}1\\4\\7\end{pmatrix},\; \begin{pmatrix}2\\5\\8\end{pmatrix}\right\}
    \]

    \textbf{Espacio nulo:} De la RREF, $x_3$ es variable libre. Con $x_3 = t$:
    \[
    x_1 = t, \quad x_2 = -2t \quad \Rightarrow \quad \mathcal{B}_{\text{nulo}} = \{(1,-2,1)\}
    \]

    \textbf{Verificación:} $\text{rango}(A) + \text{nulidad}(A) = 2 + 1 = 3$ \checkmark

    %---

    \item Para $B = \begin{pmatrix} 1 & 0 & -2 & 1 \\ 0 & 1 & 3 & -2 \\ 2 & -1 & -7 & 4 \end{pmatrix}$:
    \[
    B \xrightarrow{F_3-2F_1}
    \begin{pmatrix} 1 & 0 & -2 & 1 \\ 0 & 1 & 3 & -2 \\ 0 & -1 & -3 & 2 \end{pmatrix}
    \xrightarrow{F_3+F_2}
    \begin{pmatrix} 1 & 0 & -2 & 1 \\ 0 & 1 & 3 & -2 \\ 0 & 0 & 0 & 0 \end{pmatrix}
    \]

    \textbf{Espacio fila:} $\mathcal{B}_{\text{fila}} = \{(1,0,-2,1),\;(0,1,3,-2)\}$

    \textbf{Espacio columna:} $\mathcal{B}_{\text{col}} = \left\{\begin{pmatrix}1\\0\\2\end{pmatrix},\;\begin{pmatrix}0\\1\\-1\end{pmatrix}\right\}$

    \textbf{Espacio nulo:} Variables libres $x_3 = s$, $x_4 = t$:
    \[
    x_1 = 2s - t, \quad x_2 = -3s + 2t
    \]
    \[
    \mathcal{B}_{\text{nulo}} = \{(2,-3,1,0),\;(-1,2,0,1)\}
    \]

    \textbf{Verificación:} $2 + 2 = 4 = n$ \checkmark

    %---

    \item Para $C = \begin{pmatrix} 2 & -1 & 3 \\ 4 & -2 & 6 \end{pmatrix}$:
    \[
    C \xrightarrow{F_2-2F_1} \begin{pmatrix} 2 & -1 & 3 \\ 0 & 0 & 0 \end{pmatrix} \xrightarrow{\frac{1}{2}F_1} \begin{pmatrix} 1 & -\frac{1}{2} & \frac{3}{2} \\ 0 & 0 & 0 \end{pmatrix}
    \]

    \textbf{Espacio fila:} $\mathcal{B}_{\text{fila}} = \{(2,-1,3)\}$

    \textbf{Espacio columna:} $\mathcal{B}_{\text{col}} = \left\{\begin{pmatrix}2\\4\end{pmatrix}\right\}$

    \textbf{Espacio nulo:} Variables libres $x_2 = s$, $x_3 = t$:
    \[
    x_1 = \tfrac{1}{2}s - \tfrac{3}{2}t \quad \Rightarrow \quad \mathcal{B}_{\text{nulo}} = \{(1,2,0),\;(-3,0,2)\}
    \]

    \textbf{Verificación:} $1 + 2 = 3 = n$ \checkmark

\end{enumerate}

%-------------------------------------------------
% Solución 2: Rango
%-------------------------------------------------
\newpage
\item Rango de las matrices.

\begin{enumerate}

    \item Reduciendo a forma escalonada:
    \[
    \begin{pmatrix} 1 & 3 & -2 \\ 2 & -1 & 4 \\ -1 & 4 & -6 \end{pmatrix}
    \xrightarrow{F_2-2F_1,\;F_3+F_1}
    \begin{pmatrix} 1 & 3 & -2 \\ 0 & -7 & 8 \\ 0 & 7 & -8 \end{pmatrix}
    \xrightarrow{F_3+F_2}
    \begin{pmatrix} 1 & 3 & -2 \\ 0 & -7 & 8 \\ 0 & 0 & 0 \end{pmatrix}
    \]
    Hay 2 filas no nulas, por lo tanto $\text{rango} = 2$.

    \item Reduciendo:
    \[
    \begin{pmatrix} 0 & 1 & 2 & 3 \\ 1 & 0 & -1 & 2 \\ 2 & 1 & 0 & 7 \\ 3 & 2 & 1 & 12 \end{pmatrix}
    \longrightarrow \cdots \longrightarrow
    \begin{pmatrix} 1 & 0 & -1 & 2 \\ 0 & 1 & 2 & 3 \\ 0 & 0 & 0 & 0 \\ 0 & 0 & 0 & 0 \end{pmatrix}
    \]
    Hay 2 filas no nulas, por lo tanto $\text{rango} = 2$.

    \item Para la matriz diagonal $\text{diag}(a,b,c)$, su rango es igual al número de entradas diagonales distintas de cero:
    \begin{itemize}
        \item Si $a \neq 0$, $b \neq 0$, $c \neq 0$: $\text{rango} = 3$.
        \item Si exactamente uno es cero: $\text{rango} = 2$.
        \item Si exactamente dos son cero: $\text{rango} = 1$.
        \item Si $a = b = c = 0$: $\text{rango} = 0$.
    \end{itemize}

\end{enumerate}

%-------------------------------------------------
% Solución 3: Sistema Ax=0
%-------------------------------------------------
\newpage
\item Sistema $A\mathbf{x} = \mathbf{0}$.

\begin{enumerate}

    \item Reducción de $A$:
    \[
    \begin{pmatrix} 1 & -1 & 2 & 1 \\ 2 & -2 & 4 & 2 \\ 3 & 1 & -1 & 0 \\ 1 & 3 & -5 & -1 \end{pmatrix}
    \xrightarrow{F_2-2F_1,\;F_3-3F_1,\;F_4-F_1}
    \begin{pmatrix} 1 & -1 & 2 & 1 \\ 0 & 0 & 0 & 0 \\ 0 & 4 & -7 & -3 \\ 0 & 4 & -7 & -2 \end{pmatrix}
    \]
    \[
    \xrightarrow{F_4 - F_3}
    \begin{pmatrix} 1 & -1 & 2 & 1 \\ 0 & 4 & -7 & -3 \\ 0 & 0 & 0 & 1 \\ 0 & 0 & 0 & 0 \end{pmatrix}
    \xrightarrow{\text{RREF}}
    \begin{pmatrix} 1 & 0 & \frac{1}{4} & 0 \\ 0 & 1 & -\frac{7}{4} & 0 \\ 0 & 0 & 0 & 1 \\ 0 & 0 & 0 & 0 \end{pmatrix}
    \]

    \item \textbf{Variables pivote:} $x_1$, $x_2$, $x_4$. \quad \textbf{Variable libre:} $x_3 = t$.

    \item De la RREF: $x_1 = -\tfrac{1}{4}t$, $x_2 = \tfrac{7}{4}t$, $x_4 = 0$. Entonces:
    \[
    \mathbf{x} = t\begin{pmatrix} -\frac{1}{4} \\ \frac{7}{4} \\ 1 \\ 0 \end{pmatrix}
    = \frac{t}{4}\begin{pmatrix} -1 \\ 7 \\ 4 \\ 0 \end{pmatrix}
    \]
    Base del espacio nulo: $\mathcal{B}_{\text{nulo}} = \{(-1,7,4,0)\}$. \quad $\text{nulidad}(A) = 1$.

    \item $\text{rango}(A) = 3$. Verificación: $3 + 1 = 4 = n$ \checkmark

\end{enumerate}

%-------------------------------------------------
% Solución 4: Espacio columna y consistencia
%-------------------------------------------------
\newpage
\item Consistencia y espacio columna de $A$.

Primero reducimos $A$ para encontrar su rango y forma escalonada:
\[
\begin{pmatrix} 1 & 2 & -1 \\ 3 & 1 & 2 \\ -1 & 3 & -4 \end{pmatrix}
\xrightarrow{F_2-3F_1,\;F_3+F_1}
\begin{pmatrix} 1 & 2 & -1 \\ 0 & -5 & 5 \\ 0 & 5 & -5 \end{pmatrix}
\xrightarrow{F_3+F_2}
\begin{pmatrix} 1 & 2 & -1 \\ 0 & -5 & 5 \\ 0 & 0 & 0 \end{pmatrix}
\]

El rango de $A$ es 2. La tercera fila es siempre cero; el sistema $A\mathbf{x}=\mathbf{b}$ es consistente si y solo si la misma combinación de operaciones fila aplicada a $\mathbf{b}$ produce cero en la tercera entrada.

\begin{enumerate}
    \item Para $\mathbf{b}=(1,2,3)^T$: Augmentando y reduciendo:
    \[
    \left(\begin{array}{ccc|c} 1&2&-1&1\\3&1&2&2\\-1&3&-4&3\end{array}\right)
    \longrightarrow
    \left(\begin{array}{ccc|c} 1&2&-1&1\\0&-5&5&-1\\0&0&0&3\end{array}\right)
    \]
    La última fila representa $0=3$, contradicción. El sistema es \textbf{inconsistente}; $\mathbf{b}\notin \text{Col}(A)$.

    \item Para $\mathbf{b}=(5,8,2)^T$:
    \[
    \left(\begin{array}{ccc|c} 1&2&-1&5\\3&1&2&8\\-1&3&-4&2\end{array}\right)
    \longrightarrow
    \left(\begin{array}{ccc|c} 1&2&-1&5\\0&-5&5&-7\\0&0&0&0\end{array}\right)
    \]
    No hay contradicción. El sistema es \textbf{consistente}; $\mathbf{b}\in\text{Col}(A)$.

    \item Para $\mathbf{b}=(0,0,0)^T$: El sistema homogéneo siempre es consistente (solución trivial). $\mathbf{b}\in\text{Col}(A)$. \checkmark
\end{enumerate}

%-------------------------------------------------
% Solución 5: Cuatro espacios fundamentales
%-------------------------------------------------
\newpage
\item Cuatro espacios fundamentales de $M$.

La matriz $M$ ya está en forma escalonada reducida:
\[
M = \begin{pmatrix} 1 & 2 & 3 \\ 0 & 1 & 4 \\ 0 & 0 & 0 \end{pmatrix}
\xrightarrow{F_1-2F_2}
\begin{pmatrix} 1 & 0 & -5 \\ 0 & 1 & 4 \\ 0 & 0 & 0 \end{pmatrix}
\]

\begin{enumerate}

    \item \textbf{Espacio fila:} filas no nulas de la RREF:
    \[
    \mathcal{B}_{\text{fila}} = \{(1,0,-5),\;(0,1,4)\}, \quad \dim = 2
    \]

    \textbf{Espacio columna:} columnas 1 y 2 de $M$ original (columnas pivote):
    \[
    \mathcal{B}_{\text{col}} = \left\{\begin{pmatrix}1\\0\\0\end{pmatrix},\;\begin{pmatrix}2\\1\\0\end{pmatrix}\right\}, \quad \dim = 2
    \]

    \textbf{Espacio nulo:} variable libre $x_3 = t$, $x_1 = 5t$, $x_2 = -4t$:
    \[
    \mathcal{B}_{\text{nulo}} = \{(5,-4,1)\}, \quad \dim = 1
    \]

    \textbf{Espacio nulo izquierdo} (núcleo de $M^T$): Se reduce $M^T$:
    \[
    M^T = \begin{pmatrix}1&0&0\\2&1&0\\3&4&0\end{pmatrix}
    \longrightarrow
    \begin{pmatrix}1&0&0\\0&1&0\\0&0&0\end{pmatrix}
    \]
    Variable libre $y_3 = t$: $\;\mathcal{B}_{\text{nulo izq}} = \{(0,0,1)\}$, $\;\dim = 1$.

    \item Tabla resumen:
    \[
    \renewcommand{\arraystretch}{1.4}
    \begin{array}{|l|l|c|}
    \hline
    \textbf{Espacio} & \textbf{Base} & \textbf{Dimensión} \\
    \hline
    \text{Espacio fila} & \{(1,0,-5),\,(0,1,4)\} & 2 \\
    \text{Espacio columna} & \{(1,0,0)^T,\,(2,1,0)^T\} & 2 \\
    \text{Espacio nulo} & \{(5,-4,1)\} & 1 \\
    \text{Espacio nulo izquierdo} & \{(0,0,1)\} & 1 \\
    \hline
    \end{array}
    \]

    \item \textbf{Verificación para $M$ ($3\times 3$):}
    \[
    \text{rango}(M) + \text{nulidad}(M) = 2 + 1 = 3 = n \;\checkmark
    \]
    \textbf{Verificación para $M^T$ ($3\times 3$):}
    \[
    \text{rango}(M^T) + \text{nulidad}(M^T) = 2 + 1 = 3 \;\checkmark
    \]
    Como era de esperarse, $\text{rango}(M)=\text{rango}(M^T)=2$.

\end{enumerate}

\end{enumerate}

\end{document}
